\chapter{Mất mới thấy quý}
\markboth{Mất mới thấy quý}{Loss Reveals Value}

\section{Một người chỉ biết quý sức khỏe sau lần nằm viện.}
\begin{truyen}{Một người chỉ biết quý sức khỏe sau lần nằm viện.}{Pain clarifies priorities.}
\chuhoa{B}{ình thường anh đổi giấc ngủ lấy việc, đổi bữa cơm lấy tiện lợi. Sau vài ngày nằm viện, anh mới thấy điều quý nhất đã từng bị mình dùng như thứ luôn sẵn có.}
\textit{(He traded sleep for work and real meals for convenience. After a few days in the hospital, he saw that what was most precious had been treated as endlessly available.)}
\end{truyen}
\begin{baihoc}
Nhiều giá trị chỉ sáng lên khi ta đã chạm vào viễn cảnh mất nó.
\textit{(Many values shine clearly only when we touch the possibility of losing them.)}
\end{baihoc}
\begin{ghinhoanh}\textbf{Short English to remember:}

Pain clarifies priorities.
\end{ghinhoanh}
\begin{tuhocnhanh}
1. Em đang xem nhẹ điều quý nào? \textit{(What precious thing are you taking lightly?)}

2. Nếu mất nó, em sẽ hối tiếc điều gì? \textit{(If you lost it, what would you regret?)}
\end{tuhocnhanh}
\ngancach

\section{Người con chỉ thấy thương mẹ khi bếp nhà lạnh đi.}
\begin{truyen}{Người con chỉ thấy thương mẹ khi bếp nhà lạnh đi.}{Absence teaches quietly.}
\chuhoa{K}{hi còn ở nhà, cậu thấy mọi thứ đều quen và dễ bỏ qua. Đến lúc căn bếp không còn tiếng gọi cơm, cậu mới hiểu mình từng sống giữa rất nhiều yêu thương mà không gọi tên nổi.}
\textit{(When he still lived at home, everything felt ordinary and easy to overlook. Once the kitchen fell quiet, he realized he had lived among deep love without ever naming it.)}
\end{truyen}
\begin{baihoc}
Yêu thương quen thuộc thường bị coi nhẹ nhất cho đến khi nó biến mất khỏi nếp sống thường ngày.
\textit{(Familiar love is often undervalued most, until it disappears from daily life.)}
\end{baihoc}
\begin{ghinhoanh}\textbf{Short English to remember:}

Absence teaches quietly.
\end{ghinhoanh}
\begin{tuhocnhanh}
1. Điều yêu thương nào trong nhà em đang quá quen? \textit{(What act of love at home has become too familiar?)}

2. Em muốn biết ơn ai nhiều hơn? \textit{(Whom do you want to appreciate more?)}
\end{tuhocnhanh}
\ngancach

\section{Một người chỉ biết quý thời gian sau khi lãng phí tuổi trẻ.}
\begin{truyen}{Một người chỉ biết quý thời gian sau khi lãng phí tuổi trẻ.}{Late awareness still hurts.}
\chuhoa{N}{hìn lại những năm trôi vì trì hoãn, anh thấy đau không phải vì đã mệt, mà vì đã sống rất nhiều ngày như thể ngày nào cũng có thể dùng lại.}
\textit{(Looking back on years lost to delay, what hurt him was not fatigue, but living as if every day could be used again.)}
\end{truyen}
\begin{baihoc}
Thời gian là thứ dễ coi thường nhất lúc còn nhiều, và đắt nhất lúc nhìn lại.
\textit{(Time is easiest to underestimate when it feels abundant, and most expensive in hindsight.)}
\end{baihoc}
\begin{ghinhoanh}\textbf{Short English to remember:}

Time gets expensive in hindsight.
\end{ghinhoanh}
\begin{tuhocnhanh}
1. Em đang tiêu ngày tháng cho điều gì? \textit{(What are you spending your days on?)}

2. Điều đó có đáng với tuổi trẻ của em không? \textit{(Is it worth your youth?)}
\end{tuhocnhanh}
\ngancach

\section{Một người mất bạn rồi mới biết lời nói nặng có sức phá lớn.}
\begin{truyen}{Một người mất bạn rồi mới biết lời nói nặng có sức phá lớn.}{Words can outlive anger.}
\chuhoa{C}{ơn giận chỉ kéo dài vài phút, nhưng lời nói ra lại ở lại rất lâu. Khi tình bạn gãy đi, anh mới hiểu có những mất mát bắt đầu từ một khoảnh khắc không giữ được miệng mình.}
\textit{(Anger lasts minutes, but spoken words remain much longer. When the friendship broke, he understood that some losses begin in one moment of unguarded speech.)}
\end{truyen}
\begin{baihoc}
Lời nói không có xương nhưng có thể làm gãy những điều rất khó hàn lại.
\textit{(Words have no bones, yet they can break things that are hard to mend.)}
\end{baihoc}
\begin{ghinhoanh}\textbf{Short English to remember:}

Words can outlive anger.
\end{ghinhoanh}
\begin{tuhocnhanh}
1. Em có đang xem nhẹ sức nặng của lời mình nói? \textit{(Are you underestimating the weight of your words?)}

2. Có câu nào em nên sửa lại với ai không? \textit{(Is there something you should make right with someone?)}
\end{tuhocnhanh}
\ngancach

\section{Một người mất một cơ hội nhỏ rồi mới biết đời không đợi mãi.}
\begin{truyen}{Một người mất một cơ hội nhỏ rồi mới biết đời không đợi mãi.}{Small losses warn us early.}
\chuhoa{A}{nh từng bỏ qua một cơ hội học tập vì thấy không cần gấp. Khi những cơ hội lớn hơn cũng qua theo, anh mới hiểu đời thường cảnh báo ta bằng những lần mất nhỏ trước khi lấy đi điều lớn.}
\textit{(He once ignored a learning opportunity because it did not feel urgent. When bigger opportunities passed as well, he saw that life often warns us through small losses before larger ones.)}
\end{truyen}
\begin{baihoc}
Những mất mát nhỏ đôi khi là lời nhắc sớm để ta sống tỉnh hơn.
\textit{(Small losses can be early reminders to live more awake.)}
\end{baihoc}
\begin{ghinhoanh}\textbf{Short English to remember:}

Small losses can wake you up.
\end{ghinhoanh}
\begin{tuhocnhanh}
1. Lần mất nhỏ nào đang cảnh báo em điều gì? \textit{(What is a recent small loss warning you about?)}

2. Em có đang chịu học từ nó không? \textit{(Are you willing to learn from it?)}
\end{tuhocnhanh}
\ngancach